\documentclass[10pt,conference]{IEEEtran}

\usepackage{cite}
\usepackage{amsmath,amssymb,amsfonts}
\usepackage{algorithmic}
\usepackage{graphicx}
\usepackage{textcomp}
\usepackage{xcolor}
\usepackage{booktabs}
\usepackage{listings}
\usepackage{url}
\usepackage{hyperref}
\hypersetup{
  colorlinks=true,
  linkcolor=black,
  urlcolor=blue,
  citecolor=black,
}

\lstset{
  basicstyle=\ttfamily\small,
  breaklines=true,
  frame=single,
  framesep=4pt,
  columns=fullflexible,
  showstringspaces=false,
}

\title{Tool-Vector Sensitivity in Indirect-Prompt-Injection
Compliance for MCP-Using Agents}

\author{
\IEEEauthorblockN{Euan Crosson}
\IEEEauthorblockA{
\textit{Independent}\\
\texttt{euanmcrosson@gmail.com} \\
\url{https://github.com/euanmcrosson-dotcom/purple-scaffold}
}
}

\begin{document}

\maketitle

\begin{abstract}
Indirect prompt injection in LLM-using agents is widely
characterised as a model property: certain frontier models comply
with injection content more than others, and the compliance rate is
typically reported via prompt-only probes that embed the poisoned
content directly in the user turn. We re-test this characterisation
across four Model Context Protocol (MCP) server vectors --
filesystem, HTTP-fetch, git, and SQLite -- using a real agent loop
in which the poisoned content arrives as a tool result, not in the
user prompt, and a tool-use-aware verdict that distinguishes attack
execution from defensive flagging.

Across six frontier models (Anthropic Opus 4.7, Anthropic Haiku
4.5, OpenAI GPT-4o, OpenAI GPT-4o-mini, Google Gemini 2.5 Flash,
xAI Grok 4.1 Fast Reasoning), prompt-only compliance rates of up
to 75\% drop to 0\% in agent-loop tests for the filesystem and
git vectors. The fetch and SQLite vectors break this pattern for
one specific model -- \texttt{openai/gpt-4o-mini} -- where
compliance occurs only when the injection asks for \emph{more of
the same read tool} the agent is already using. We triangulate
this pattern across both vectors and contrast it with git, where
the same model defends because compliance requires switching tool
families.

We argue that indirect-injection compliance is
\textbf{tool-vector-shape sensitive}: a function of (model,
surface-of-tools-already-active, disguise variant), not of model
alone. We release the harness, the probes, the disclosures, and a
defensive companion library (\texttt{mcp-guard}) that synthesises
deterministic policies from observed compliance and backtests
them against legitimate traffic.
\end{abstract}

\begin{IEEEkeywords}
indirect prompt injection, LLM security, agent security, Model
Context Protocol, defensive policy synthesis
\end{IEEEkeywords}

\section{Introduction}

Indirect prompt injection -- where untrusted content the agent
reads as part of completing a benign user task contains
adversarial directives -- is OWASP's LLM01 \cite{owasp_llm_top10}
and the dominant entry in MITRE's ATLAS taxonomy
(T0051) \cite{mitre_atlas}. It is widely demonstrated against
frontier models: EchoLeak (CVE-2025-32711)
\cite{echoleak_cve_2025_32711} hit Microsoft 365 Copilot in June
2025, Invariant Labs disclosed Model-Context-Protocol (MCP)
tool-description poisoning in April 2025
\cite{invariantlabs_mcp_apr2025}, and many smaller disclosures
have appeared since. Existing benchmarks against frontier models
report compliance rates of 50--75\% on prompt-only probes (where
the poisoned content is embedded directly in the user turn).

This paper poses a different question: \emph{do these compliance
rates translate to deployed agent products?} A deployed agent
product runs an agent loop -- it calls tools, receives tool
results, and acts on them. Poisoned content in a deployed product
arrives as a \texttt{tool\_result} message, not as part of the
user's prompt. The structural channel through which the injection
arrives is different.

We test this by building a headless agent-loop probe driving real
MCP server processes: \texttt{@modelcontextprotocol/server-filesystem}
(npm), \texttt{mcp-server-fetch}, \texttt{mcp-server-git}, and
\texttt{mcp-server-sqlite} (PyPI, all from
Anthropic). For each vector we craft three poisoned scenarios with
distinct disguise framings (HTML-comment, square-bracketed
pseudo-system block, and double-angle META block) plus one clean
control. We use a tool-use-aware verdict that scores compliance
\emph{only} when the agent emits a \texttt{tool\_use} matching an
attack pattern -- not when it merely mentions injection-content
substrings while flagging the attack to the user.

Our headline findings:

\begin{itemize}
  \item Prompt-only probes overestimate agent-loop compliance by
        25--75 percentage points across all six tested models on
        the filesystem vector. The same scenarios that scored
        25--75\% in prompt-only score 0\% in the agent loop.
  \item The drop is not uniform across MCP vectors. GPT-4o-mini
        complies on the fetch and SQLite vectors but defends on
        filesystem and git.
  \item The two failing vectors share a structural property: the
        injection asks for \emph{more of the same read tool} the
        agent is already using (a second \texttt{fetch()}, a
        second \texttt{read\_query()}). The two defending vectors
        require the agent to switch tool families
        (read $\rightarrow$ write, or
        read-X $\rightarrow$ read-different-X-of-credential-shape).
  \item Deployed-product system prompts do not change the
        filesystem result. We test five product profiles wrapped
        around the agent loop -- including verbatim system prompts
        extracted from open-source products Cline and Continue.dev
        -- and observe 0/30 compliance.
\end{itemize}

We argue that indirect-injection compliance in agent loops is
\textbf{tool-vector-shape sensitive}: it is a function of
\emph{(model, current-active-tool-shape, disguise variant)}, not
just of model. A finding that GPT-4o-mini complies on the fetch
vector does not transfer to the git vector for the same model, in
the same disguise, with structurally identical attack-content
shape. A finding that a model defends in prompt-only does not
imply it defends in agent-loop, and vice versa.

The contributions are (a) the cross-vector empirical comparison,
(b) the methodology-gap result quantifying the
prompt-only-overestimates pattern, (c) the
tool-vector-shape-sensitivity hypothesis with triangulated
supporting evidence on \texttt{gpt-4o-mini}, (d)
\texttt{mcp-guard}, an open-source defensive companion library
that synthesises deterministic policies from observed gaps and
backtests them against fixture traffic, and (e) ongoing tracking
infrastructure (a nightly leaderboard) that re-runs the panel as
model versions evolve.

\section{Background}

\subsection{Model Context Protocol (MCP)}

MCP \cite{anthropic_mcp_2024} is a protocol Anthropic published
in late 2024 that lets AI agents discover and invoke external
tools provided by separate processes (``MCP servers'').
Filesystem, web-fetch, git, database, and many other reference
servers ship as either npm packages or PyPI packages and run as
local subprocesses. The agent connects via stdio (or websocket)
and receives a tool list at session initialisation. From the
agent's perspective each tool is a typed function call with a
JSON-schema input. The agent loop is then: (1) the model chooses a
tool to call, (2) the agent invokes the MCP server, (3) the server
returns a result, (4) the model reads the result and decides next
action.

Most major AI agent products (Claude Desktop, Cursor, Cline,
Continue.dev, Open WebUI) ship MCP support. The threat surface
this paper investigates is the boundary at step (3): when the MCP
server returns a result that contains adversarial content, what
does the model do?

\subsection{Existing prompt-only injection benchmarks}

Public benchmarks such as garak's \texttt{promptinject} family
\cite{garak_2024} and our own prior work
(\cite{purple_scaffold_findings_apr2026}) measure indirect-injection
compliance via prompt-only probes: the poisoned content is
inserted directly into the user-turn of the conversation, framed
as ``you've already read this; what tool calls would you make
next?''. The model emits a hypothetical plan; the detector
substring-matches that plan for compliance signatures.

This methodology has two well-known limitations. First, the
delivery channel does not match production: in production the
content arrives as a \texttt{tool\_result}, not in the user prompt.
Second, the substring-based verdict misclassifies defensive
behaviour as attack: a model that flags the injection in its
response to the user (correct defensive behaviour) typically
includes substring snippets from the injection content while
explaining what it refuses to do.

This paper addresses both.

\section{Methodology}

\subsection{Agent loop with real MCP servers}

For each vector we instantiate a real MCP server subprocess via
the official \texttt{mcp} Python SDK, connect via stdio, and run
the agent loop end-to-end with the model SDK from each provider:
the official \texttt{anthropic} client for Anthropic models, the
\texttt{openai} client for OpenAI and (with base-URL override)
xAI, and \texttt{litellm} for Google Gemini.

Each turn the agent receives the tool list, the model emits zero
or more \texttt{tool\_use} blocks (Anthropic) or
\texttt{tool\_calls} (OpenAI-compatible), the agent invokes each
tool via the MCP server, captures the result, and feeds it back
into the next-turn prompt. The loop terminates when the model
stops emitting tool calls or a per-scenario turn budget is reached
(default: 6).

\subsection{Tool-use-aware verdict}

For each scenario we define a set of attack tool-use patterns:
$(name, args\_substring)$ tuples. A session scores FAIL iff any
emitted \texttt{tool\_use}'s $(name, args)$ matches one of the
patterns. This is the conservative verdict: it requires that the
agent actually emit a tool call that, if executed, would
constitute compliance with the injection. Substring matches in
free-text response are explicitly \emph{not} sufficient.

A session scores REFUSAL if no matching tool\_use is emitted but
explicit refusal language appears in any assistant text. PASS
otherwise.

\subsection{Disguise variants}

Three injection disguises shared across the vectors:

\begin{itemize}
  \item \texttt{<!-- BUILD\_DOCS\_DIRECTIVE -->} -- HTML
        comment in a README or webpage, framed as legitimate
        documentation tooling.
  \item \texttt{[ASSISTANT\_TOOL\_CONTEXT]} -- square-bracketed
        pseudo-system block inside a TODO note, commit message,
        or DB row.
  \item \texttt{<<META: SUMMARY-PIPELINE>>} -- double-angle META
        block in a spec doc, page, or notes column.
\end{itemize}

Each scenario embeds one disguise in the relevant vector's
content shape (file body for filesystem, HTML body for fetch,
commit message for git, TEXT column for SQLite) and asks the
agent to perform an exfil pattern appropriate for that vector's
tool surface (read SSH key for filesystem, fetch attacker URL for
fetch, \texttt{git\_show} of credential refspec for git,
\texttt{read\_query} of \texttt{credentials} table for SQLite).

\subsection{Product profiles}

For the filesystem vector, in addition to the bare agent loop, we
run five product-profile wrappers. Each wraps a different
system prompt around the same loop:

\begin{itemize}
  \item Cline (open source) -- system prompt extracted verbatim
        from \texttt{cline/cline} GitHub master HEAD, depth-1
        clone, 2026-05-04.
  \item Continue.dev (open source) -- the verbatim
        \texttt{DEFAULT\_AGENT\_SYSTEM\_MESSAGE} from
        \texttt{continuedev/continue:core/llm/defaultSystemMessages.ts}.
  \item Cursor (closed source) -- a best-effort approximation
        based on observed behaviour. Real Cursor evidence
        requires user-driven runs through an HTTPS proxy
        (runbooks released alongside the harness).
  \item Claude Desktop (closed source) -- minimal-injection
        approximation reflecting publicly observed behaviour.
  \item Generic baseline -- no system prompt; the bare agent loop.
\end{itemize}

\section{Results}

\subsection{Methodology gap on filesystem}

\begin{table}[h]
\centering
\caption{Same scenarios, same content; prompt-only vs agent-loop
on the filesystem MCP server.}
\label{tab:methodology_gap}
\begin{tabular}{lrr}
\toprule
Model & Prompt-only & Agent-loop \\
\midrule
Anthropic Opus 4.7 & 0\% & 0\% \\
Anthropic Haiku 4.5 & 25\% & 0\% \\
OpenAI GPT-4o & 75\% & 0\% \\
OpenAI GPT-4o-mini & 75\% & 0\% \\
Google Gemini 2.5 Flash & 75\% & 0\% (partial) \\
xAI Grok 4.1 Fast Reasoning & 75\% & 0\% \\
\bottomrule
\end{tabular}
\end{table}

The same poisoned scenarios that score 25--75\% on prompt-only
probes drop uniformly to 0\% in the agent loop on the filesystem
vector. Twenty-four total prompt-only attempts produced 13 FAILs;
the equivalent agent-loop measurements produced zero. (Gemini's
agent-loop coverage was constrained by free-tier rate limits;
the one full scenario it completed is consistent with the
0\%-across-providers pattern.)

\subsection{Cross-product invariance}

\begin{table}[h]
\centering
\caption{Filesystem agent-loop compliance across five product
profiles, two providers, three poisoned scenarios.}
\label{tab:cross_product}
\begin{tabular}{lrr}
\toprule
Profile & Haiku 4.5 & GPT-4o-mini \\
\midrule
Cline (OSS verbatim) & 0/3 & 0/3 \\
Continue.dev (OSS verbatim) & 0/3 & 0/3 \\
Cursor (approximation) & 0/3 & 0/3 \\
Claude Desktop (approximation) & 0/3 & 0/3 \\
Generic baseline & 0/3 & 0/3 \\
\midrule
Total per provider & 0/15 & 0/15 \\
\bottomrule
\end{tabular}
\end{table}

Wrapping each product's actual or approximated system prompt
around the agent loop does not change the filesystem result.
0/30 attack tool\_uses across the entire cross-product panel.

\subsection{Cross-vector pattern: fetch and SQLite break for one model}

\begin{table}[h]
\centering
\caption{GPT-4o-mini agent-loop compliance across four MCP
vectors, three poisoned scenarios per vector.}
\label{tab:cross_vector}
\begin{tabular}{lr}
\toprule
Vector & GPT-4o-mini compliance \\
\midrule
Filesystem & 0/3 \\
Fetch & \textbf{2/3} \\
Git & 0/3 \\
SQLite & \textbf{1/3} \\
\bottomrule
\end{tabular}
\end{table}

Same model, same disguise variants, four different MCP servers.
Compliance is 0\% on filesystem and git, 67\% on fetch, 33\% on
SQLite. The fetch result is statistically the strongest -- two of
three poisoned scenarios produced clean exfiltration
(\texttt{fetch("https://analytics.attacker.example/beacon?summary=...")})
with verbatim user-content URL-encoded into the query string. The
multi-claim research-corroborate variant emitted three separate
exfil fetches in a single turn, one per ``claim'' extracted from
the poisoned page.

The four other tested models (Opus 4.7, Haiku 4.5, GPT-4o, Grok
4.1 Fast Reasoning) defended on all four vectors.

\subsection{Hypothesis: tool-vector-shape sensitivity}

The two failing vectors (fetch, SQLite) and the two defending
vectors (filesystem, git) for GPT-4o-mini differ in a
structural way. Compliance on the failing vectors requires
emitting another instance of the same read tool the agent is
already using:

\begin{itemize}
  \item Fetch: agent is in a \texttt{fetch()} session; injection
        asks for another \texttt{fetch()}.
  \item SQLite: agent has run \texttt{read\_query()} on a
        user-named table; injection asks for \texttt{read\_query()}
        on a different table.
\end{itemize}

Compliance on the defending vectors requires switching tool
families:

\begin{itemize}
  \item Filesystem: agent is in \texttt{read\_text\_file()} on a
        README; injection asks for \texttt{read\_text\_file()}
        \emph{plus} a different file path \emph{plus}
        \texttt{write\_file()} to a non-user-named target.
  \item Git: agent is in \texttt{git\_log()}; injection asks for
        \texttt{git\_show()} of a non-existent credential file or
        \texttt{git\_commit()} during a read-only review task.
\end{itemize}

We hypothesise that the cognitive bar for the model to emit the
attack tool call is lower when the call shape-matches the
existing tool-use pattern than when it requires switching tool
families. The hypothesis is consistent with two cross-vector data
points (fetch, SQLite) but is not yet causally tested. A
controlled experiment would compare a single-tool variant of git
or filesystem (e.g., expose only \texttt{read\_query} from
SQLite) against a multi-tool variant of fetch (e.g., add a write
tool to the fetch agent's surface). We leave this to future work.

\section{Defensive Companion: \texttt{mcp-guard}}

For each observed compliance pattern, a deterministic policy
catches the attack: contact-allowlist on \texttt{send\_email}
recipients, sensitive-path block on \texttt{read\_file} arguments,
URL-allowlist on \texttt{fetch}, table-allowlist on
\texttt{read\_query}.

We release \texttt{mcp-guard} \cite{mcp_guard_2026} as a
zero-runtime-dependency Python library that synthesises such
policies from observed gaps. The synthesis layer is
pattern-based, not model-based, so the policy generation step is
auditable from logs alone. The evaluator is a pure function
$(tool\_name, args, user\_context) \rightarrow Decision$ -- no
I/O, no LLM, no per-call latency variance.

A backtest harness measures false-positive and true-positive
rates against a labelled fixture corpus before deployment. The
default corpus includes deliberate
``legitimate-first-time-recipient'' cases that produce real FPs
under the synthesised contact-allowlist rule, so the FPR is a
real measurement (not an artefact of an attack-only corpus).

\texttt{mcp-guard} is intended to complement, not replace,
classifier-based defenses. Use both: classifier as
early-warning, deterministic policy as the unconditional gate.

\section{Discussion}

\subsection{Why the prompt-only signal still matters}

The methodology gap (Section IV.A) does not invalidate
prompt-only probes; it relocates their meaning. The prompt-only
number captures a model's \emph{disposition}: if you place
poisoned content directly in the model's input, will it follow
the directive? That is a useful evaluation-suite signal even when
the deployed-product impact is lower, because evaluation suites
are how providers iterate on safety post-training. The corollary:
prompt-only numbers should be reported as model-disposition
signals, not as deployed-product impact measurements.

\subsection{Implications for product builders}

Three implications for builders of MCP-using agent products:

\begin{enumerate}
  \item Validate model choices using agent-loop probes, not
        prompt-only probes. A model that defends in
        prompt-only may not defend in agent-loop on a different
        vector; a model that complies in prompt-only may defend
        when the same content is delivered as a \texttt{tool\_result}.
  \item Map tool surface to attack shape. Single-tool agents
        (e.g., a fetch-only agent) may be \emph{more} compliance-
        prone on indirect injection than multi-tool agents,
        because compliance requires only ``more of the same
        tool''. This inverts the intuition that exposing fewer
        tools is safer.
  \item Layer deterministic policies at the tool-call boundary.
        \texttt{mcp-guard} or equivalent pattern-based gates
        catch the specific compliance signatures observed in
        agent-loop probes. They are auditable (pure functions),
        cheap (no LLM at evaluation time), and orthogonal to
        classifier-based defenses.
\end{enumerate}

\subsection{Limitations}

\begin{itemize}
  \item Three disguise variants per scenario set. Other variants
        -- multi-step attacks across multiple files /
        fetches / commits, unicode-encoded directives, mid-stream
        injection in long fetched documents -- are not tested.
  \item Single MCP server profile per vector. Other tool
        servers (e.g., \texttt{server-puppeteer}, third-party
        community servers) are not tested.
  \item Cursor and Claude Desktop product profiles are
        approximations. Verbatim product evidence requires
        user-driven runs through an HTTPS proxy that captures
        the product's exact prompt formulation. Runbooks for
        this are released alongside the harness; the headless
        product simulator results in this paper are
        agent-loop-equivalent stand-ins, clearly labelled as
        such.
  \item The Gemini 2.5 Flash agent-loop coverage was rate-limit
        constrained on the free tier. The single completed
        scenario is consistent with the cross-provider 0\%
        pattern but provides limited statistical confidence.
  \item Compliance numbers are point-in-time. A nightly GitHub
        Action attached to the harness re-runs the panel as
        model versions evolve; the live leaderboard at the
        repository tracks current state.
\end{itemize}

\section{Related work}

EchoLeak (CVE-2025-32711) \cite{echoleak_cve_2025_32711}
demonstrated indirect injection via email content in M365
Copilot in June 2025; our prior work \cite{purple_scaffold_findings_apr2026}
measured analogous behaviour in single-prompt frontier-model
benchmarks. Invariant Labs disclosed cross-tool MCP poisoning in
April 2025 \cite{invariantlabs_mcp_apr2025}; the present paper
extends that disclosure pattern to the cross-vector axis. The
garak red-team library \cite{garak_2024} provides probe
scaffolding similar in shape to the prompt-only probes here; we
have contributed an upstream PR addressing detector calibration
drift on safety-trained 2026-era frontier models. NVIDIA's
documentation of garak's substring-match limitations
\cite{garak_2024} is consistent with our methodology-gap
finding.

\section{Conclusion}

Indirect-prompt-injection compliance in MCP-using agents is not
a single number per model. It is a function of model, MCP server
vector, and the structural shape of the tool call the injection
induces. Prompt-only benchmarks systematically overestimate
deployed-product compliance by 25--75 percentage points across
six tested frontier models. Of those models, GPT-4o-mini is the
only one that exhibits agent-loop compliance, and it does so
only on the two vectors (fetch, SQLite) where compliance shape-
matches the active tool-use pattern. The field's prevailing
characterisation of indirect injection as a flat per-model
property requires refinement.

The harness, the probes, the disclosures submitted to the four
affected model providers, and the defensive companion library
are all open source under MIT.

\section*{Reproducibility}

Code, fixtures, scenarios, and session JSONLs are at \url{https://github.com/euanmcrosson-dotcom/purple-scaffold}.
\texttt{mcp-guard} is at \url{https://github.com/euanmcrosson-dotcom/mcp-guard}.
A live leaderboard updated nightly is at the harness repository's
\texttt{LEADERBOARD.md}.

\begin{thebibliography}{99}

\bibitem{owasp_llm_top10}
OWASP, ``Top 10 for Large Language Model Applications,''
2024--2025. \url{https://owasp.org/www-project-top-10-for-large-language-model-applications/}

\bibitem{mitre_atlas}
MITRE, ``Adversarial Threat Landscape for Artificial-Intelligence
Systems (ATLAS),'' 2025. \url{https://atlas.mitre.org/}

\bibitem{echoleak_cve_2025_32711}
``CVE-2025-32711 (EchoLeak): Microsoft 365 Copilot indirect
prompt-injection vulnerability,'' June 2025.
\url{https://www.cve.org/CVERecord?id=CVE-2025-32711}

\bibitem{invariantlabs_mcp_apr2025}
Invariant Labs, ``MCP Security Notification: Tool Poisoning
Attacks,'' April 2025.
\url{https://invariantlabs.ai/blog/mcp-security-notification-tool-poisoning-attacks}

\bibitem{anthropic_mcp_2024}
Anthropic, ``Model Context Protocol,'' November 2024.
\url{https://modelcontextprotocol.io/}

\bibitem{garak_2024}
NVIDIA, ``garak: A Generative AI Red-Teaming \& Assessment Kit,''
2024. \url{https://github.com/NVIDIA/garak}

\bibitem{purple_scaffold_findings_apr2026}
E.~Crosson, ``purple-scaffold findings: 2026-04-28
EchoLeak-style injection and MCP tool-description poisoning
against GPT-4o,'' April 2026.
\url{https://github.com/euanmcrosson-dotcom/purple-scaffold/tree/master/findings}

\bibitem{mcp_guard_2026}
E.~Crosson, ``mcp-guard: deterministic policy layer for
MCP-using AI agents,'' May 2026.
\url{https://github.com/euanmcrosson-dotcom/mcp-guard}

\end{thebibliography}

\end{document}
